\chapter{Генерація, валідація і верифікація}

Комплексна інфраструктура рівня ERP вимагає бездоганного теоретичного фундаменту для забезпечення безвідмовної роботи алгоритмів. В рамках архітектурного підходу (відомого як \emph{Volume V: Verification} або \emph{Computational Layer}) ми визначаємо набір математичних та інженерних дисциплін. Програмне забезпечення повинно бути не лише емпірично протестованим, але й строго доведеним за допомогою формальних методів. Розділ закладає теоретичний фундамент, оглядаючи сучасний стан формальної верифікації та визначаючи місце пропонованої уніфікованої системи формальних мов у цій галузі.

Формальна верифікація як дисципліна бере початок із робіт Ніколаса де Брейна, чия система AUTOMATH (1967) стала першим механізмом комп’ютеризованого доведення теорем у фіброваній моделі. Хенк Барендрегт розвинув цю ідею, створивши лямбда-куб — класифікацію типових систем, де найвища вершина (Calculus of Constructions, CoC) уможливила вищі формальні мови. Ми використовуємо CoC як базис для теоретичного ядра, доповнюючи його вищими мовами для поглинання машинною верифікацією усієї доступної математики підприємства.

\section{Формальна верифікація та валідація}

Для унеможливлення помилок на виробництві застосовуються різні методи формальної верифікації. Формальна верифікація — доказ, або заперечення відповідності системи у відношенні до певної формальної специфікації або характеристики, із використанням формальних методів математики.

Згідно з міжнародними нормами (IEEE, ANSI)\footnote{IEEE Std 1012-2016 --- V\&V Software verification and validation} та у відповідності до вимог Європейського Аерокосмічного Агенства\footnote{ESA PSS-05-10 1-1 1995 -- Guide to software verification and validation}, процес валідації включає в себе перегляд (code review), тестування (модульне, інтеграційне, властивостей), перевірку моделей, аудит — увесь комплекс необхідний для доведення, що продукт відповідає вимогам висунутим при розробці.

Огляд сучасних засобів формальної верифікації показує, що найкращі результати досягаються тоді, коли мова програмування сама слугує інструментом математичного доведення. Класифікація цих підходів виділяє окрему \emph{спектральну категорію мов формальної верифікації} (Dependently Typed Systems). Процес розподіляється на валідацію (чи правильну систему ми будуємо?) та верифікацію (чи правильно ми її побудували?).

\section{Формальна специфікація}

Для спрощення процесу верифікації та валідації застосовується математична техніка формалізації постановки задачі --- формальна специфікація. Це математична модель, створена для опису систем, визначення їх основних властивостей, та інструментарій для перевірки цих систем.

Існують два фундаментальні підходи:
1) Алгебраїчний підхід, де система описується в термінах операцій та відношень між ними (аналітичний метод);
2) Модельно-орієнтований підхід, де модель створена конструктивними побудовами на базі теорії множин, а системні операції визначаються тим, як вони змінюють стан системи (синтетичний метод).

\subsection{Програмне забезпечення та логіка}

Найбільш стандартизована та прийнята в області формальної верификації --- це нотація Z\footnote{ISO/IEC 13568:2002 --- Z formal specification notation} (Spivey, 1992), приклад модельно-орієнтованої мови, названої на честь Ернеста Цермело.

Інша відома мова формальної специфікації як стандарт для моделювання розподілених систем, таких як телефонні мережі та протоколи, це LOTOS\footnote{ISO 8807:1989 --- LOTOS} як приклад алгебраїчного підходу. Ця мова побудована на темпоральних логіках та поведінках залежниих від спостережень. Інші мови специфікацій, які можна відзначити тут --- це TLA+, CSP, CCS (Milner, 1971), Actor Model, BPMN, тощо.

\subsection{Математичні компоненти}
Перші системи комп'ютерної алгебри розроблялися ще під PDP-6: MATHLAB, MACSYMA, SCRATCHPAD, REDUCE, SACLIB, MUMATH. Сучасні системи включають AXIOM, MAGMA, MUPAD, GAP.

\section{Формальні методи верифікації}

Можна виділити три підходи до верифікації:
1) спеціалізовані верифікатори моделей, або системи моделювання (VST, NuPRL, TLA+, Twelf, SystemVerilog);
2) алгебраїчні мови для синтетичних моделей і глибокого вбудовування (Coq, Agda, Lean, F*, cubicaltt, RedPRL);
3) системи автоматичного доведення теорем і синтезу програм (HOL/Isabelle, ACL2).

\subsection{Алгебраїчні мови та System F}
Вступ у теорію мов програмування неможливий без розуміння фундаменту типізації, побудованого на поліморфному лямбда-численні \textbf{System F}. Завдяки System F, програми автоматично гарантують відсутність цілих класів помилок періоду виконання. Компілятор розв'язує рівняння рівності типів через механізми уніфікації, працюючи у безпечній та математично обґрунтованій категорії.

\section{Моделі процесів}

Для моделювання складних бізнес-процесів (BPMN) та високонавантажених мережевих протоколів ми спираємося на математичну формалізацію, здатну автоматизувати генерування тестів. Для цього розроблено спеціалізовану мову специфікацій \textbf{Zen Crypted Dharma DSL (Z.180)}, яка слугує виконавчим доповненням до міжнародного стандарту \textbf{ITU-T Z.120 Message Sequence Chart (MSC)}. 

Мова Z.180 абстрагує транспортний рівень та криптографічні кодування, природним чином маплячись на конструкти Z.120: сесії (Sessions) стають інстансами (lifelines), дії (Actions) — вихідними/вхідними повідомленнями, а експектації (Expectations/Predicates) — умовами перевірки. Завдяки цьому забезпечується строга верифікація міжшарових інваріантів (консистентність курсорів читання, механізми ABAC, модерації) без втручання в імутабельну історію подій. Цей підхід безпосередньо сумісний з TTCN-3 для генерації тест-сьютів і формальною семантикою процес-алгебр з Annex B. Процес генерації інтеграційних шлюзів та SDK на основі цих специфікацій стає рутинною механічною трансформацією з AST у клієнти Swift, Kotlin, Erlang, C\# та JavaScript.

\section{Формальні мови та середовища виконання}

Усі середовища виконання можна умовно розділити на два класи:
1) інтерпретатори нетипізованого або просто типізованого лямбда числення (Erlang/BEAM, V8, HotSpot, Kx, PyPy);
2) безпосередня генерація інструкцій процессора і лінкування (OCaml, Rust, Haskell, Pony).

Найбільш цікаві цільові платформи для виконання программ які побудовані на основі формальних доведень для нас є OCaml (основна мова екстракту Coq), Rust (відсутність сміттєзбірника), Erlang (неблокована семантика $\pi$-числення) та Pony.

\subsection{Формальні інтерпретатори та екстракція}
Рантайм може будуватися як інтерпретатор нетипізованої системи. У рамках розвитку проєкту використовується PTS тайпчекер Henk (на базі Erlang/OCaml) у якості проміжної мови для повної нормалізації лямбда термів. Тотальна програма виступає способом лінкування з підсистемою вводу-виводу віртуальної машини Erlang, поєднуючи функціональну довершеність (CoC) з реальною неблокованою роботою.

\section{Базова схема підприємства ERP/1}

Генерація, валідація та верифікація зливаються в єдиний обчислювальний конвеєр для бази підприємства. Та встановлено, що усі системи доведення і середовища виконання є носіями мовних елементів, які згідно теорії типів Мартіна-Льофа можна кластеризувати так:
\begin{itemize}
    \item $O_\lambda$ --- нетипизоване $\lambda$-числення Чорча;
    \item $O_\pi$ --- числення процесів, CCS, CSP або $\pi$-числення Мілнера;
    \item $O_\mu$ --- тензорне числення та векторизація;
    \item $O_\Pi$ --- числення конструкцій (функціональна повнота);
    \item $O_\Sigma$ --- числення контекстів (контекстуальна повнота);
    \item $O_=$ --- теорія типів Мартіна-Льофа (логіка);
    \item $O_W$ --- числення індуктивних конструкцій (матіндукція);
    \item $O_I$ --- гомотопічна система типів (формальна математика).
\end{itemize}

Усі бізнес-правила (BPMN/BPE), протоколи міжсервісної взаємодії та політики доступу (ABAC) проходять етап символьного виконання, тестування за моделлю MSC та типового розв'язання. Лише код, який успішно формує доведення власної консистентності у термінах цих вищих мов програмування, допускається до середовища виконання підприємства.
